\chapter{氧、硫、氮、磷：生命与环境的骨架}

\begin{kaichang}
找一盒火柴（没有就去厨房看一眼燃气灶），做一件你做过无数次、却未必细想过的事：划一根火柴。

“刺啦”一声，火柴头先冒出一小股白烟，随即腾起火焰，还飘来一点刺鼻的气味。你可能不知道，这一秒钟里登场的角色几乎全是这一章的主角：火柴盒侧面那条粗糙的红褐色涂层里，藏着\textbf{磷}；火柴头里混着\textbf{硫}，以及一种受热会放出\textbf{氧}的白色粉末（氯酸钾）。摩擦生的热点燃磷，磷点燃硫，硫引燃木梗，而火柴头自己释放的氧气让火烧得更旺。

现在回答两个问题。第一：火明明可以在空气里烧，火柴头为什么还要自带“氧气包”？第二：空气里最多的气体根本不是氧气——将近五分之四是氮气。整根火柴都泡在氮气里，为什么氮气从头到尾袖手旁观，既不帮忙也不捣乱？

氧、硫、氮、磷，这四种元素为什么偏偏凑在一起，既撑起了你的骨骼和能量供应，又搅动着酸雨和湖泊？先猜，再往下看。
\end{kaichang}

\section{现象：一种火上浇油，一种冷眼旁观}

先只看能观察到的现象。

关于氧气，你可能在课堂上演示里见过：把一根带火星的木条伸进盛有氧气的集气瓶，木条“腾”地复燃，烧得又亮又猛。蜡烛在纯氧中燃烧，火焰比在空气里旺得多；细铁丝在空气里烧不起来，在纯氧里却能火星四射。反过来，把燃着的蜡烛扣进密闭的玻璃罩，火焰会慢慢变弱、熄灭——罩子里的“助燃成分”被用掉了。

关于氮气，现象正好相反。鼓鼓的薯片袋里充的就是氮气：它把氧气挤走，让薯片不被氧化哈喇，又当“气垫”防碎。液氮则是另一副面孔：温度低到约 $-196\,^{\circ}\mathrm{C}$，能让花瓣一敲就碎，却被用来冷冻保存细胞和种子——因为它冷归冷，几乎不和任何东西反应。

铁生锈是慢动作的氧气现象：一颗铁钉放在潮湿角落里，几周就蒙上红锈。这个反应慢到几乎看不见，但它确实在消耗空气里的某种成分。下面的实验能让你“看见”这份消耗。

\begin{shiyan}
\textbf{看得见的“空气被吃掉一块”}

\textbf{材料}：一小团钢丝绒（五金店有售）、一个窄口玻璃瓶、一个浅盘、水、几滴白醋。

\textbf{步骤}：钢丝绒蘸一点白醋（洗去表面防锈油），塞进玻璃瓶底部；盘里倒水，把玻璃瓶倒扣进盘里，让瓶内外水面相平，用笔标记水位。放在窗边，每天观察，持续三到五天。

\textbf{观察点}：钢丝绒逐渐变成红棕色；瓶内水面缓缓上升。到最后，上升的水大约占瓶内空气体积的五分之一。

\textbf{安全提示}：全程不需要火，也不要尝试用任何方法“制造纯氧”来加速。钢丝绒边缘锋利，做完洗手；白醋别溅进眼睛。

上升的这几厘米水，是氧气离开空气、进入铁锈留下的“空位”。粗略地，你在家里就量出了空气中氧气的份额——约五分之一。
\end{shiyan}

\section{粒子：氧分子牵两只手，氮分子牵三只}

为什么同样是空气的主要成分，一个活泼、一个懒惰？把镜头推进到分子层。

回忆第 10 章的周期表街区图：氧和硫住在第 16 族，最外层 6 个电子；氮和磷住在第 15 族，最外层 5 个电子。对它们来说，“再拉来一两个电子把外层填得更满”往往对应更低的能量（注意：这是能量规律，不是原子的愿望）。所以这四家的元素都有较强的非金属脾气，爱跟别的原子共用或夺取电子——这是它们能成为生命骨架的前提。

空气里的氧是两个氧原子组成的分子 \chem{O_2}，两个原子之间共用两对电子，像两只手紧紧相握（第 19 章会细讲共价键）。氮气也是双原子分子 \chem{N_2}，但两个氮原子之间共用\textbf{三对}电子——三只手交握。别小看这多出来的一只手：拆开 \chem{N_2} 中两个氮原子，需要约 \ud{945}{kJ\,mol^{-1}} 的能量，是拆开 \chem{O_2}（约 \ud{498}{kJ\,mol^{-1}}）的近两倍，在常见双原子分子里几乎是最结实的。

于是现象层的对比有了粒子层的解释：氮气不是“不想”反应，而是拆开它的“启动费”太贵。常温下分子碰撞的能量大多付不起这笔启动费，反应就被挡在门外——这是\textbf{速度}问题，不是方向问题（第 27、29 章会把“能不能”和“快不快”彻底分开讲）。氧气启动费低得多，所以木头、铁丝、你体内的葡萄糖，都付得起。

现在让符号正式登场。木头的主要成分简化成碳来写，燃烧就是：
\[\chem{C} + \chem{O_2} \rightarrow \chem{CO_2}\]
铁生锈慢得多，但也是同类过程：铁原子把电子交给氧，生成铁的氧化物。你每分每秒进行的呼吸，账本上写的也是“有机物 $+$ 氧气 $\rightarrow$ 二氧化碳 $+$ 水”——和燃烧同方向的账本，只是生命把它拆成几十步、慢慢来（第 54 章细讲）。

\section{氧气不是燃料，它是“收电子的一方”}

这里必须停一下，拆一个流传极广的说法。

\begin{wuqu}
“氧气能燃烧”“氧气是燃料”——这是错的，而且这个误解有真实的危险。燃料是被氧化的那一方（木头、汽油、氢气），氧气是\textbf{接受电子}的那一方，叫氧化剂。氧气自己从来不烧：你把纯氧灌进气球，它一点就着不了——因为没有燃料。真正危险的是反过来：在纯氧环境里，平时温和的东西（油脂、衣物、头发）都会变成猛烈燃料，所以医院氧气瓶附近严禁油污和明火。顺便说，用“得氧失氧”定义氧化还原只是历史起点，本质上是电子的转移，第 33 章会换成电子语言重新讲。
\end{wuqu}

所以燃烧需要三样：燃料、氧气、足够高的启动温度，缺一样就灭。盖灭蜡烛是断了氧气，浇水灭火主要是降了温度。这也回答了开场第一个问题的一半：火柴头自带“氧气包”（氯酸钾受热放出 \chem{O_2}），是要在火柴头那一小团地方瞬间制造出\textbf{远比空气浓}的氧，确保火一定能着——空气里那 21\% 的氧，对小火苗来说经常不够“立竿见影”。

\begin{qiaoliang}
\textbf{一口空气里有多少氧气？}干燥空气按体积算：氮约 78\%，氧约 21\%，剩下约 1\% 主要是氩和二氧化碳。你的教室如果长 9 米、宽 7 米、高 3 米，体积约 $189\ \mathrm{m^3}$，其中氧气约 $189\times 21\% \approx 40\ \mathrm{m^3}$。一个静坐的中学生每分钟大约消耗 $0.25\ \mathrm{L}$ 氧气，一天约 $0.36\ \mathrm{m^3}$——教室里的氧理论上够一个人用一百多天。那为什么会觉得闷？因为让你难受的往往不是缺氧，而是\textbf{二氧化碳升高}：呼出气里二氧化碳从约 0.04\% 升到约 4\%，它在血液中的积累才是“闷”的主要信号。感觉和测量常常不是一回事，这正是要用仪器而不是用直觉的原因。
\end{qiaoliang}

\section{臭氧：同一个分子，两种名声}

氧还有一种三个原子的形态：臭氧 \chem{O_3}。它结构不如 \chem{O_2} 稳定，性情更急躁，是更强的氧化剂。有趣的是，臭氧的名声完全取决于它待在哪里。

在距离地面 20 到 30 千米的高空，臭氧形成薄薄一层。如果把这层臭氧压到地面常温常压，厚度只有约 3 毫米——就这么薄的一层，却吸收了阳光中绝大部分对生物杀伤最强的紫外线。没有它，DNA 会被阳光直接打断化学键，陆地生命很难存在。在这个高度上，臭氧是地球撑开的遮阳伞。

而在你呼吸的地面，臭氧是污染物。夏日阳光强烈时，汽车尾气和工业排放的氮氧化物、挥发性有机物在光照下发生一连串反应，生成臭氧和其他刺激性分子，形成光化学烟雾。近地面臭氧刺激眼睛和呼吸道，对哮喘患者尤其危险；它还会损伤植物叶片。同一种分子，高空是保护伞，地面是刺激物——\textbf{位置决定角色}。

这也是理解一切环境化学的钥匙之一：谈一种物质“好”还是“坏”，必须先说清它在哪里、浓度多高。消毒柜和印刷机附近也会产生臭氧，闻起来有股淡淡的“雷雨天味道”——雷电确实能把少量 \chem{O_2} 打成 \chem{O_3}，但那点浓度远不足以“净化室内空气”，臭氧机在高浓度下对人同样有害。

上世纪后半叶，人类还发现了一个意外：广泛用于制冷剂的氟利昂释放到高空后，在紫外线作用下放出氯原子，一个氯原子能反复催化分解上万个臭氧分子，在南极上空撕出了“臭氧洞”。这段历史留到后面讲催化剂和自由基时再细拆，这里只记住一点：一个看似惰性的分子，换个高度，就能变成拆伞的剪刀。

\section{氮：拆不动的三键，和一场喂饱一半人类的工业}

回到开场第二个问题：火柴泡在氮气里，氮气为什么旁观？现在答案只剩一层窗户纸：\chem{N_2} 的三键太结实，划火柴那点温度付不起拆开它的启动费。

但“难拆”不等于“拆不动”。雷电瞬间的高温能撕开三键，让氮和氧化合成氮氧化物，随雨水落进土壤——这是天然氮肥。豆科植物的根瘤里住着一类细菌，体内的酶能在常温下把 \chem{N_2} 一步步改造成氨——这是生物固氮，第 54 章前后还会再遇到这些微观化工厂。问题在于：自然固氮的速度，远远追不上人口增长后农田对氮的需求。

为什么氮这么要紧？因为蛋白质、DNA 里都少不了氮原子，而植物只能吸收土壤里“现成的”含氮离子（硝酸盐、铵盐），不能直接吃空气里的 \chem{N_2}——就像一个守着米仓却拿不到钥匙的人。空气里每吨氮气都在头顶飘，庄稼却在地里缺氮。谁能把空气里的氮“固定”成植物能用的形态，谁就握住了粮食产量的阀门。

\begin{zhengju}
19 世纪末，欧洲农业越来越依赖从智利海运来的硝石矿。1898 年，英国科学家威廉·克鲁克斯在一次著名演讲中警告：硝石照这个速度消耗，几十年内就会见底，小麦产量将撑不住人口，“以面包为主食的文明”会挨饿。问题被清晰地提了出来：怎样把空气里用不完的 \chem{N_2} 变成氨？

当时的化学知识已经能算出一笔矛盾的账：合成氨反应放热，温度越低，平衡越偏向氨这边；可温度低了，\chem{N_2} 的三键又拆得太慢，慢到没有实用价值。著名化学家能斯特据此认为这条路在工业上走不通。德国化学家哈伯不服，他做了关键的事：\textbf{把反应搬到高压下}。加压让气体分子更密集，平衡向分子数更少的产物一侧移动，产率被硬生生抬高。1909 年，哈伯的实验装置在约 200 个大气压、$500\,^{\circ}\mathrm{C}$ 上下、含铁催化剂的条件下，让氮气和氢气源源不断地变成氨，产率数据推翻了能斯特的悲观估计——能斯特本人验证后公开承认了错误。

但实验室装置不等于工厂：耐高压的钢罐、廉价氢气、能长期工作的催化剂，每一项都是工程难关。博施带领巴斯夫公司的团队花了四年多，试了两万多种配方寻找催化剂，才在 1913 年建成第一座合成氨工厂。

这个故事的尾声并不轻松。合成氨让化肥产量摆脱了矿藏的限制，今天全世界用这个方法固定的氮，养活了大约一半人口；但同一套装置在两次世界大战中也被用来生产炸药的原料硝酸。哈伯晚年因犹太身份被自己效忠的国家驱逐。一项技术从实验室走进世界，带来的从来不只是它最初承诺的那一面——这个主题在本书第三册还会反复出现。
\end{zhengju}

合成氨的账本写出来只有一行：
\[\chem{N_2} + \chem{3H_2} \rightleftharpoons \chem{2NH_3}\]
可逆符号提醒你：这是一个能来回跑的平衡（第 31 章细讲），工业上一边不断移走氨、一边把没反应的气体压回去循环，硬是把“懒气体”变成了流水线上的原料。

\begin{qiaoliang}
\textbf{你的身体里有多少“哈伯的氮”？}今天全球合成氨工业每年固定约 1.2 亿吨氮。摊到 80 亿人头上，人均约 15 千克氮——相当于每人每年分到约 90 千克硝酸铵化肥里的氮量。再看人体：一个体重 60 千克的人，氮约占体重的 3\%，也就是约 1.8 千克，几乎全在蛋白质和 DNA 里。研究者估算，当今人类体内大约一半的氮原子，追根溯源都来自合成氨工厂固定的氮。换句话说，你身上大约 0.9 千克的氮，曾经在一座高压钢罐里被从 \chem{N_2} 中拆出来。元素不认路牌，它只在“空气—化肥—庄稼—你”之间流动。
\end{qiaoliang}

\section{磷：骨头里的石头、细胞里的“能量零钱”}

氮的楼下邻居是磷。同属第 15 族，最外层也是 5 个电子，但磷是固体，性格远没有氮那么“宅”。

单质磷有两副著名的面孔。白磷分子是四个磷原子搭成的小四面体 \chem{P_4}，四个原子挤在小角度里，键被掰得很不舒服，能量偏高，所以白磷暴躁到在空气里 $40\,^{\circ}\mathrm{C}$ 左右就能自燃，必须泡在水里保存（这是实验室里的事，你绝不应该在家接触它）。红磷的结构则连成更大的网络，张力小得多，性格温和，要加热到约 $240\,^{\circ}\mathrm{C}$ 才着火——所以它被安全地刷在火柴盒侧面，靠摩擦的瞬间高温才登场。同一种元素，仅仅因为原子连接方式不同，脾气天差地别：这个主题（同素异形体）在第 36 章讲碳时还会大放异彩。

生命对磷的用法要沉稳得多。你的骨骼和牙齿约八成多的无机成分是羟基磷酸钙，一种含钙、磷酸根和氢氧根的矿物——骨头本质上是“石头框架加胶原钢筋”的复合材料。每个细胞膜的骨架磷脂分子里有磷酸根；DNA 的螺旋“扶梯扶手”正是糖和磷酸根交替连成的长链（第 58 章细讲）。而磷最出名的一个角色，藏在三个字母里：ATP。

ATP（三磷酸腺苷）常被叫作细胞的“能量货币”，但流行的讲法“能量藏在高能磷酸键里，断开就释放”是\textbf{错的}——请记住一条铁律：\textbf{断开任何化学键都要吸收能量，能量只在形成新键时放出}（第 27 章会专门讲）。ATP 水解真正发生的事是：反应
\[\chem{ATP} + \chem{H_2O} \rightarrow \chem{ADP} + \chem{Pi}\]
中，产物（ADP 和磷酸根）比反应物稳定得多——ATP 分子里挤在一起的几个负电荷被分开，排斥减小，产物还能以更多方式分散到水中。断开旧键要花钱，形成新键和分散电荷赚回来更多，\textbf{净赚}的差额才是细胞可用的能量。ATP 更像一枚“找得开的零钱”：不是因为它体内藏着火药，而是因为把它花出去（水解）这笔交易整体划算。

磷在环境里的故事则是另一副面孔：\textbf{富营养化}。磷在地壳里不算多，又几乎不进入大气，所以在大多数湖泊和河流中，磷是水藻最缺的那一口“饭”——其他营养再多，没有磷，藻也暴不起来；磷一来，藻就疯长。含磷洗衣粉、农田流失的磷肥，曾经让大量水体绿成油漆：藻类暴发后死亡，细菌分解藻尸时大量耗氧，鱼和其他动物缺氧而死。这就是许多国家禁售含磷洗衣粉的原因。氮和磷怎样在岩石、土壤、生物和水体之间循环流动，是理解生态系统的关键账本，第 16 章会完整展开。

\section{硫：臭鸡蛋、硫酸厂和变酸的雨}

氧的楼下邻居是硫。同族，最外层同样 6 个电子，但比氧多一层电子楼，性格跟着微妙地变了：抢电子的劲头比氧弱一档，却因此能干一些氧干不了的活。

单质硫是黄色的脆性固体，常温下最稳定的形态是八个硫原子连成一顶“皇冠”环 \chem{S_8}。硫和生命的关系可以从气味说起：臭鸡蛋、坏掉的卷心菜、火山口的刺鼻味，都是含硫小分子在报警——臭鸡蛋的主角是硫化氢 \chem{H_2S}，由细菌分解蛋白质里的含硫氨基酸产生。你的头发和角蛋白富含硫原子之间形成的“二硫键”，烫发药水的原理就是拆开再重接这些硫桥。顺便提醒：\chem{H_2S} 高浓度时有毒，而且它会麻痹嗅觉——闻不到不代表安全，所以下水道、沼气池附近的气味警告必须当真，这是只能由专业人员处理的事。

硫在工业上的地位，可以用一句话概括：一个国家的硫酸产量，常被当作工业化程度的体温计。化肥、电池、冶金、化纤、药品，几乎行行用硫酸。现代制硫酸的主线是：硫或含硫矿石先在空气中烧成二氧化硫 \chem{SO_2}，再在催化剂（常用五氧化二钒）表面被氧气进一步氧化成三氧化硫 \chem{SO_3}，最后吸收进水里：
\[\chem{S} + \chem{O_2} \rightarrow \chem{SO_2}\]
\[\chem{2SO_2} + \chem{O_2} \rightleftharpoons \chem{2SO_3}\]
\[\chem{SO_3} + \chem{H_2O} \rightarrow \chem{H_2SO_4}\]
三步里第二步最磨人：它也需要“付启动费”，催化剂的作用是给反应换一条更便宜的路径——注意，催化剂只改变速度，不改变反应的方向和最终的收支（第 30 章专门讲）。

但 \chem{SO_2} 一旦从烟囱直接排进大气，麻烦就来了。煤和石油里天生含硫，燃烧时硫变成 \chem{SO_2} 随烟气上天，在空气和云水里被逐步氧化、溶解，最终变成稀硫酸落回地面；汽车尾气里的氮氧化物走类似的路变成硝酸。这样的雨就是酸雨。这里有个容易忽略的事实：\textbf{正常雨水本来就微酸}，pH 约 5.6，因为空气中的 \chem{CO_2} 溶进雨水会生成少量碳酸；酸雨指的是 pH 明显低于这个背景值的降水。酸雨腐蚀石灰岩雕像和建筑（大理石主要成分是碳酸钙，酸能把它溶解），让湖泊变酸、鱼类死亡，还会把土壤里的铝溶出来伤害树根。治理思路也很“化学”：电厂烟气用石灰石浆液“洗澡”，让 \chem{SO_2} 和碳酸钙反应被固定下来；或者干脆烧低硫煤。问题从元素出发，答案也得回到元素。

\section{这套说法在哪里会失灵}

本章反复用了“活泼”“懒惰”“暴躁”“温和”这类词。它们是好用的速记，但都有明确的边界，越界就会误导你。

第一，活泼与懒惰都是\textbf{条件}下的结论，不是元素的永久人格。氮气在常温下袖手旁观，在雷电的高温里、在哈伯的高压罐里、在根瘤菌的酶旁边，照样乖乖反应。氧气在纯氧瓶里让铁丝剧烈燃烧，在常温空气里让铁钉生锈却要等上几周——“能不能反应”和“多快反应”是两本账，分别由能量规则和速率规则掌管（第 27、29 章）。把“惰性”当成“永远不反应”，第 13 章讲稀有气体时你已经见过一次反例了。

第二，“某元素有益”或“某元素有害”的说法没有科学意义。臭氧在高空护命、在地面伤肺；磷是骨骼和 ATP 的命脉，进了湖泊就是灾难；硫化氢微量时是气味信号、还是体内的一种信号分子，高浓度时致命；就连氧气，早产儿吸入过量也会损伤视网膜。物质的角色由\textbf{位置、剂量、化学形态}三者共同决定——这条原则请带进你以后看到的每一条“养生”或“排毒”广告里。

第三，本章的收支账都是“典型条件下”的近似。真实体系里常有多个反应竞争：富营养化的水体里，限制藻类的不一定总是磷，有些湖泊和近海是氮在当“瓶颈”；酸雨的成因里硫酸和硝酸的比重也因地区能源结构而异。科学模型给出的不是万能答案，而是“先看哪个变量”的优先级。

\section{再看那根火柴}

现在回到开场，把那一秒钟重放一遍，这次带上粒子层的眼镜。

你的手指给火柴头一个摩擦，摩擦生热，局部温度瞬间升高。火柴盒侧面的红磷在高温下有一小部分转化为更活泼的形态并被点燃；火星溅到火柴头上，头里的氯酸钾受热分解，原地释放出高浓度的氧气——这是自带的“氧气包”，比空气里 21\% 的氧迅猛得多。浓氧扑向火柴头里的硫：
\[\chem{S} + \chem{O_2} \rightarrow \chem{SO_2}\]
硫燃烧放出的热又点燃了木梗，你闻到的那股刺鼻味，正是 \chem{SO_2}。火焰沿着木梗蔓延，是因为木质纤维（主要是纤维素，碳、氢、氧的化合物）不断被加热到能付得起“启动费”，与氧气持续结账。

而从头到尾占空气五分之四的氮气，确实就在火焰旁边流过。它不是不想参与，是三键的启动费太贵，火柴的温度付不起大头——它只是被加热、膨胀、流走，最多在火焰最热的缝隙里生成一丁点氮氧化物。

一根火柴，四种元素，各按各的键能和条件出牌。所谓化学，不过是把这副牌桌看清楚。

\section{顺着链条继续走}

这一章我们第一次让“元素街区”走出了周期表，进入空气、骨骼、农田和烟囱。链条没有停：氧、硫、氮、磷的故事会在三个方向上继续。第一，哈伯用的铁催化剂为什么能替反应“另辟蹊径”？铁、钒、锰这些过渡金属凭什么在表面上撮合分子？那是第 15 章要逛的街区。第二，氮和磷在自然界的来路和去处——它们怎样在岩石、海洋、生物之间循环，人类又怎样把这些循环改得面目全非——第 16 章会画出完整的物质流地图。第三，ATP 这枚“零钱”在细胞里怎样赚来又怎样花掉，要等到第二册讲呼吸作用的第 54 章才算总账。

\begin{tiaozhan}
\textbf{为什么合成氨工厂偏要选在“两头不讨好”的条件？}（跳过不影响主线。）合成氨反应放热，且反应后气体分子数减少（4 个分子变 2 个）。用热力学看：低温让平衡更偏向氨，高压也让平衡偏向氨——那为什么不把温度降到室温、压力开到极限？因为动力学拦路：温度越低，拆开 \chem{N_2} 三键的启动费越难凑齐，反应慢到工厂一天出不了一罐氨。高压则受限于钢材强度和压缩能耗，压力翻倍，设备成本和事故风险远不止翻倍。于是工业上选了一个折中：约 200 个大气压、$400\sim 500\,^{\circ}\mathrm{C}$，外加铁催化剂提速，再把没反应完的 \chem{N_2} 和 \chem{H_2} 分离出来循环进料。每一趟循环只有一部分气体变成氨，但循环足够多次，总转化率照样很高。平衡、速率、成本、安全四个变量互相拉扯，最后的“最优解”是一个工程妥协，而不是某一条定律的最大值——真实世界里的化学，几乎总是这样。
\end{tiaozhan}

\section*{练一练}

\begin{sikaoti}
\item 用“共用几对电子”的画法，分别画出 \chem{O_2} 和 \chem{N_2} 分子，并在旁边画一条“拆开分子所需能量”的示意柱图。用这张图向同学解释：为什么薯片袋充氮气而不是充氧气来保鲜？
\item 画一张“臭氧位置—角色”对照图：横轴是高度（地面 / 高空 20 千米以上），标出臭氧在两个位置分别做什么、对人分别意味着什么。再用一句话总结：为什么“臭氧是好是坏”这个问题本身问错了？
\item 哈伯的高压实验推翻了能斯特的悲观估计。请说清楚：能斯特的推理错在哪里（漏掉了哪个变量）？这个故事对“理论计算和实验谁说了算”有什么启示？
\item 画一条酸雨的物质流箭头图：从煤里的硫原子出发，经过烟囱、大气、云水，到湖泊里的鱼。在每个箭头上注明发生了什么变化（是物理移动还是化学反应，生成了什么）。
\item 一间教室体积约 $189\ \mathrm{m^3}$，氧气约占 21\%。班里有 45 名同学，每人每分钟消耗约 $0.25\ \mathrm{L}$ 氧气。假设教室完全不通风，估算同学们多久会消耗掉教室里 10\% 的氧气？再想一想：现实中“闷”的感觉会比你的计算来得早得多，为什么？（提示：回顾本章关于二氧化碳的那段。）
\item 有人宣传“臭氧机可以给卧室消毒、增加氧气，有益健康”。用本章的知识指出这段话里至少两处问题，并说明判断依据。
\end{sikaoti}
